\subsection{Convex subsets of a numerical space}

Let E be the n-dimensional numerical space \(^{*}\) (or, more generally, a topological vector space over R)*. For every pair \((x,y)\) of points of E, one calls segment (resp. open segment) with endpoints x and y (cf. EVT, II, p. 7) the set of points of E of the form \(tx + (1-t)y\) , for \(t \in [0,1]\) (resp. for \(t \in ]0,1[\) ). Let A be a subset of E. One says that the set A is convex if for every pair \((x,y)\) of points of A and every \(t \in I\) , the point \(tx + (1-t)y\) belongs to A.

*A convex subset is path-connected.*

Lemma 1. --- Let E be the n-dimensional numerical space and let A be a convex and compact subset of E of which 0 is an interior point. For every \(x \in E\), denote by \(p_{\mathrm{A}}(x)\) the lower bound in \(\overline{R}\) of the set of real numbers \(t > 0\) such that \(x \in tA\).

The map \(p_{A}\) is finite, continuous, and satisfies the following properties:

\begin{enumerate}
\def\labelenumi{(\roman{enumi})}
\item
  For every \(x \in E\) such that \(x \neq 0\) , we have \(p_{\mathrm{A}}(x) > 0\) ;
\item
  For every \(s \in R_{+}\) and every \(x \in E\) , we have \(p_{\mathrm{A}}(sx) = sp_{\mathrm{A}}(x)\) .
\item
  For all \(x\) and \(y \in \mathrm{E}\) , we have \(p_{\mathrm{A}}(x + y) \leqslant p_{\mathrm{A}}(x) + p_{\mathrm{A}}(y)\) .
\item
  For a point \(x\) of E to belong to A, it is necessary and sufficient that \(p_{\mathrm{A}}(x) \leqslant 1\) .
\end{enumerate}

For \(x \in E\), denote by \(\|x\|\) the Euclidean norm of x (TG, VI, p. 7). Since A is compact, there exists a real number \(M > 0\) such that every point x of A satisfies \(\|x\| \leqslant M\). Since 0 is an interior point of A, there exists a real number \(m > 0\) such that every point x of E such that \(\|x\| \leqslant m\) belongs to A. Consequently, one has the relation \(\|x\|/\mathrm{M} \leqslant p_{\mathrm{A}}(x) \leqslant \|x\|/m\), for every \(x \in E\). In particular, \(p_{\mathrm{A}}(0) = 0\) and \(p_{\mathrm{A}}(x) \neq 0\) if \(x \neq 0\).

Assertions (ii) and (iv) follow immediately from the definition of the map \(p_{A}\) .

Let \(x\) and \(y\) be points of E. Let \(x'\) and \(y'\) be points of A such that \(x = p_{\mathrm{A}}(x)x'\) and \(y = p_{\mathrm{A}}(y)y'\) . If \(x\) and \(y\) are not both zero, \(p_{\mathrm{A}}(x) + p_{\mathrm{A}}(y) > 0\) and one has

\[
\begin{array}{l} x + y = p _ {\mathrm{A}} (x) x ^ {\prime} + p _ {\mathrm{A}} (y) y ^ {\prime} \\ \qquad = (p _ {\mathrm{A}} (x) + p _ {\mathrm{A}} (y)) \left(\frac {p _ {\mathrm{A}} (x)}{p _ {\mathrm{A}} (x) + p _ {\mathrm{A}} (y)} x ^ {\prime} + \frac {p _ {\mathrm{A}} (y)}{p _ {\mathrm{A}} (x) + p _ {\mathrm{A}} (y)} y ^ {\prime}\right). \end{array}
\]

Since A is convex, this shows that \(x + y\) belongs to \((p_{\mathrm{A}}(x) + p_{\mathrm{A}}(y))\mathrm{A}\) , whence \(p_{\mathrm{A}}(x + y) \leqslant p_{\mathrm{A}}(x) + p_{\mathrm{A}}(y)\) . If x = y = 0, this inequality is still satisfied, because \(p_{\mathrm{A}}(0) = 0\) . This proves assertion (iii).

Applying this inequality to \(x + y\) and -y, one deduces that, for every pair \((x, y)\) of points of E, we have

\[
| p _ {\mathrm{A}} (x + y) - p _ {\mathrm{A}} (x) | \leqslant \max (p _ {\mathrm{A}} (y), p _ {\mathrm{A}} (- y)) \leqslant m ^ {- 1} \| y \|.
\]

This proves that the map \(p_{A}\) is continuous, whence the lemma.

The map \(p_{A}\) is called the gauge of the convex set A.

PROPOSITION 2. --- Let E be the n-dimensional numerical space and let A be a convex and compact subset of E of which 0 is an interior point. There exists a unique bijection u of E onto itself satisfying the following three properties:

\begin{enumerate}
\def\labelenumi{(\roman{enumi})}
\item
  For every \(x \in E\) and every \(t \in R_{+}\) , \(u(tx) = tu(x)\) ;
\item
  For every \(x \in E\) , there exists \(\lambda \in R_{+}\) such that \(u(x) = \lambda x\) ;
\item
  We have \(u(\mathrm{A}) = \mathbf{B}_n\)
\end{enumerate}

The map u is a homeomorphism and induces, by passage to subspaces, a homeomorphism of A onto \(B_{n}\) , a homeomorphism of the interior of A onto the interior of \(B_{n}\) , and a homeomorphism of the boundary of A in E onto the sphere \(S_{n-1}\) .

Let \(p_{A}\) be the gauge of the convex set A. Let \(x \in E\) and let \(t \in R_{+}\); for \(x\) to belong to \(tA\), it is necessary and sufficient that \(p_{\mathrm{A}}(x) \leqslant t\). Since A is compact, there exists a real number M such that \(\|x\| \leqslant M\) for all \(x \in A\).

Let \(u\) be a map satisfying the conditions of the proposition. We have \(u(0) = 0\). Let \(x \in \mathrm{E}\setminus\{0\}\) and let \(\lambda \in \mathbf{R}_+\) such that \(u(x) = \lambda x\). For all \(t \in \mathbf{R}_+\) such that \(tx \in \mathrm{A}\), we have \(u(tx) = t\lambda x\). Since \(u\) is injective, \(\lambda \neq 0\). Since \(u(\mathrm{A})\) is contained in \(\mathbf{B}_n\), we also have \(\lambda \leqslant p_{\mathrm{A}}(x)/\|x\|\). Set \(z = x/\|x\|\); it is a point of \(\mathbf{S}_{n-1}\). For \(z\) to have a preimage in A under \(u\), it is necessary and sufficient that the point \((\lambda \|x\|)^{-1}x\) belong to A, that is, \(\lambda \|x\| \geqslant p_{\mathrm{A}}(x)\), i.e. \(\lambda \geqslant p_{\mathrm{A}}(x)/\|x\|\). This implies the uniqueness of such a map \(u\).

Let us then denote by \(u\) the map of E into itself which sends 0 to 0 and \(x\) onto \((p_{\mathrm{A}}(x) / \| x\|)x\), for every \(x \in \mathrm{E}\setminus\{0\}\).

By lemma 1, \(u\) is continuous at every point of \(\mathrm{E}\setminus\{0\}\). We have \(\| u(x)\| = p_{\mathrm{A}}(x)\) for all \(x\in \mathrm{E}\) and \(p_{\mathrm{A}}(x)\to 0\) when \(x\to 0\) (loc. cit.); consequently, \(u\) is continuous at 0. Hence, \(u\) is continuous.

The only preimage of 0 under u is 0. Let \(y \in E\setminus\{0\}\). For an element x of E to satisfy \(u(x) = y\), it is necessary and sufficient that \(x = (\|y\| / p_{\mathrm{A}}(y))y\). This entails that u is a continuous bijection of E onto itself. Its inverse is the map \(v: y \mapsto (\|y\| / p_{\mathrm{A}}(y))y\). Since \(p_{A}\) is continuous and vanishes only at 0, the map v is continuous at every point of \(E\setminus\{0\}\); the inequality \(p_{\mathrm{A}}(y) \geqslant \|y\| / \mathrm{M}\) entails that \(\|v(y)\| \leqslant \mathrm{M}\|y\|\) for all \(y \in E\). It follows that v is continuous. Hence, the map u is a homeomorphism of E onto itself. Since the relations \(p_{\mathrm{A}}(x) \leqslant 1\) and \(x \in A\) are equivalent, we also have \(u(\mathrm{A}) = \mathbf{B}_{n}\). The proposition follows.

Example. --- The set \([0,1]^{n}\) is a convex, compact subset of \(R^{n}\) with nonempty interior. It follows from proposition 2 of I, p. 123 that it is homeomorphic to the closed Euclidean ball. More precisely, for every point x of the interior \(]0,1[^{n}\) and every point b of the open Euclidean ball, there exists a homeomorphism of \([0,1]^{n}\) onto \(B_{n}\) which sends x to b and induces by passage to subspaces a homeomorphism of \(]0,1[^{n}\) onto the open Euclidean ball, as well as a homeomorphism of the boundary of \([0,1]^{n}\) onto the sphere \(S_{n-1}\).

Consequently, every convex, compact subset of \(R^{n}\) with nonempty interior is homeomorphic to a parallelepiped (loc. cit.).

